% arara: pdflatex
% arara: bibtex
% arara: makeglossaries
% arara: pdflatex
% arara: remover: { patterns: [ '*.aux', '*.xdy', '*.toc', '*.out', '*.gl*', '*.ac*', '*.alg', '*.blg', '*.bbl' ] }
%% Преамбула TeX-файла

\documentclass[utf8, 14pt, bold]{G7-32}
\sloppy

% Настройки стиля ГОСТ 7-32
% Для начала определяем, хотим мы или нет, чтобы рисунки и таблицы нумеровались в пределах раздела, или нам нужна сквозная нумерация.
\EqInChapter % формулы будут нумероваться в пределах раздела
\TableInChapter % таблицы будут нумероваться в пределах раздела
\PicInChapter % рисунки будут нумероваться в пределах раздела

% Добавляем гипертекстовое оглавление в PDF
\usepackage[
bookmarks=true, colorlinks=true, unicode=true,
urlcolor=black,linkcolor=black, anchorcolor=black,
citecolor=black, menucolor=black, filecolor=black,
]{hyperref}

\AfterHyperrefFix

\usepackage{microtype}% полезный пакет для микротипографии, увы под xelatex мало чего умеет, но под pdflatex хорошо улучшает читаемость

% Тире могут быть невидимы в Adobe Reader
\ifInvisibleDashes
\MakeDashesBold
\fi

\usepackage{graphicx}   % Пакет для включения рисунков

% С такими оно полями оно работает по-умолчанию:
% \RequirePackage[left=20mm,right=10mm,top=20mm,bottom=20mm,headsep=0pt,includefoot]{geometry}
% Если вас тошнит от поля в 10мм --- увеличивайте до 20-ти, ну и про переплёт не забывайте:
\geometry{right=10mm}
\geometry{left=30mm}
\geometry{bottom=20mm}
\geometry{ignorefoot}% считать от нижней границы текста


% Пакет Tikz
\usepackage{tikz}
\usetikzlibrary{arrows,positioning,shadows}

% Произвольная нумерация списков.
\usepackage{enumerate}

% ячейки в несколько строчек
\usepackage{multirow}

% itemize внутри tabular
\usepackage{paralist,array}

%\setlength{\parskip}{1ex plus0.5ex minus0.5ex} % разрыв между абзацами
\setlength{\parskip}{1ex} % разрыв между абзацами
\usepackage{blindtext}

% Центрирование подписей к плавающим окружениям
%\usepackage[justification=centering]{caption}

\usepackage{newfloat}
\DeclareFloatingEnvironment[
placement={!ht},
name=Equation
]{eqndescNoIndent}
\edef\fixEqndesc{\noexpand\setlength{\noexpand\parindent}{\the\parindent}\noexpand\setlength{\noexpand\parskip}{\the\parskip}}
\newenvironment{eqndesc}[1][!ht]{%
    \begin{eqndescNoIndent}[#1]%
\fixEqndesc%
}
{\end{eqndescNoIndent}}

\makeatletter
\def\thenomenclature{%
  \if@intoc\addcontentsline{toc}{section}{\nomname}\fi%
\nompreamble
\list{}{%
\labelwidth\nom@tempdim
\leftmargin\labelwidth
\advance\leftmargin\labelsep
\itemsep\nomitemsep
\let\makelabel\nomlabel}
}
\makeatother

\usepackage{array, booktabs, hhline, tabularx, indentfirst, calc, textcase, xstring}
\usepackage{stackrel}
\usepackage{ulem}
\usepackage{adjustbox}

\usepackage{titlesec}


\ifPDFTeX
	% 8 Листинги

\usepackage{listings}

% Значения по умолчанию
\lstset{
  basicstyle= \footnotesize,
  breakatwhitespace=true,% разрыв строк только на whitespacce
  breaklines=true,       % переносить длинные строки
%   captionpos=b,          % подписи снизу -- вроде не надо
  inputencoding=koi8-r,
  numbers=left,          % нумерация слева
  numberstyle=\footnotesize,
  showspaces=false,      % показывать пробелы подчеркиваниями -- идиотизм 70-х годов
  showstringspaces=false,
  showtabs=false,        % и табы тоже
  stepnumber=1,
  tabsize=4,              % кому нужны табы по 8 символов?
  frame=single
}

% Стиль для псевдокода: строчки обычно короткие, поэтому размер шрифта побольше
\lstdefinestyle{pseudocode}{
  basicstyle=\small,
  keywordstyle=\color{black}\bfseries\underbar,
  language=Pseudocode,
  numberstyle=\footnotesize,
  commentstyle=\footnotesize\it
}

% Стиль для обычного кода: маленький шрифт
\lstdefinestyle{realcode}{
  basicstyle=\scriptsize,
  numberstyle=\footnotesize
}

% Стиль для коротких кусков обычного кода: средний шрифт
\lstdefinestyle{simplecode}{
  basicstyle=\footnotesize,
  numberstyle=\footnotesize
}

% Стиль для BNF
\lstdefinestyle{grammar}{
  basicstyle=\footnotesize,
  numberstyle=\footnotesize,
  stringstyle=\bfseries\ttfamily,
  language=BNF
}

% Определим свой язык для написания псевдокодов на основе Python
\lstdefinelanguage[]{Pseudocode}[]{Python}{
  morekeywords={each,empty,wait,do},% ключевые слова добавлять сюда
  morecomment=[s]{\{}{\}},% комменты {а-ля Pascal} смотрятся нагляднее
  literate=% а сюда добавлять операторы, которые хотите отображать как мат. символы
    {->}{\ensuremath{$\rightarrow$}~}2%
    {<-}{\ensuremath{$\leftarrow$}~}2%
    {:=}{\ensuremath{$\leftarrow$}~}2%
    {<--}{\ensuremath{$\Longleftarrow$}~}2%
}[keywords,comments]

% Свой язык для задания грамматик в BNF
\lstdefinelanguage[]{BNF}[]{}{
  morekeywords={},
  morecomment=[s]{@}{@},
  morestring=[b]",%
  literate=%
    {->}{\ensuremath{$\rightarrow$}~}2%
    {*}{\ensuremath{$^*$}~}2%
    {+}{\ensuremath{$^+$}~}2%
    {|}{\ensuremath{$|$}~}2%
}[keywords,comments,strings]

% Подписи к листингам на русском языке.
\renewcommand\lstlistingname{Листинг}
\renewcommand\lstlistlistingname{Листинги}

\else
	\usepackage{local-minted}
\fi

% Любимые команды
\newcommand{\Code}[1]{\textbf{#1}}

\newcommand\arraybslash{\let\\\@arraycr}
\newcolumntype{L}[1]{>{\raggedright\let\newline\\\arraybackslash\hspace{0pt}}m{#1}}
\newcolumntype{C}[1]{>{\centering\let\newline\\\arraybackslash\hspace{0pt}}m{#1}}
\newcolumntype{R}[1]{>{\raggedleft\let\newline\\\arraybackslash\hspace{0pt}}m{#1}}
\renewcommand{\tabularxcolumn}[1]{m{#1}}
\newcolumntype{Y}{>{\centering\arraybackslash}X}


\newcommand\RPZ[1]{
	\def\RPZTitle{#1}
}

\newcommand\RPZTheme[1]{
	\def\RPZTitleTheme{#1}
}

\newcommand\RPZGroup[1]{
	\def\RPZTitleGroup{#1}
}

\newcommand\RPZStudent[1]{
	\def\RPZTitleStudent{#1}
}

\newcommand\RPZTeacher[1]{
	\def\RPZTitleTeacher{#1}
}
\def\SignatureLength{\maxof{\widthof{\RPZTitleTeacher}}{\widthof{\RPZTitleStudent}}}



\newcommand\Signature[1]{
\noindent
\hfill
\begin{minipage}{7cm}
	\begin{flushright}
		\null\hfill 
			$\stackrel
				[\text{\rmfamily\fontsize{9pt}{9pt}\selectfont (Подпись, дата)}]
				{}
				{\uline{\parbox{2.5cm}{\text{\rmfamily\fontsize{12pt}{12pt}\selectfont \null}}\rule[35pt]{0pt}{0pt}}}$ 
			\hfill 
			$\stackrel
				[\text{\rmfamily\fontsize{9pt}{9pt}\selectfont (И.О. Фамилия)}]
				{}
				{\uline{\parbox{\SignLen ex + 10 pt}{\centering\rmfamily\fontsize{12pt}{12pt}\selectfont #1} \rule[35pt]{0pt}{0pt}}}$
	\end{flushright}
\end{minipage}
}

\renewcommand{\thesection}{\arabic{section}}

\begin{document}
\frontmatter

\RPZ{к научно-исследовательской работе}
\RPZTheme{\null\hfill \bigskip Современные методы  \hfill\null \\ \\ \null\hfill \bigskip алгебраического криптоанализа \hfill\null \\ \\}%ugly
\RPZGroup{ИУ8-94}
\RPZStudent{А.Д. Полухина}
\RPZTeacher{П.Г. Ключарев}


\StrLen{\RPZTitleStudent}[\StudentLen]
\StrLen{\RPZTitleTeacher}[\TeacherLen]
\ifthenelse{\StudentLen > \TeacherLen}{\StrLen{\RPZTitleStudent}[\SignLen]}{\StrLen{\RPZTitleTeacher}[\SignLen]}

\pagestyle{empty}
\noindent
\begin{tabularx}{\textwidth}{L{2cm} C{\textwidth - 2.7cm}}
    \includegraphics[scale=1]{inc/img/bmstu.jpg} &
    {\centering\rmfamily\bfseries\fontsize{11pt}{11pt}\selectfont
    Министерство науки и высшего образования Российской Федерации \\
    Федеральное государственное бюджетное образовательное учреждение \\
    высшего образования \\
    <<Московский государственный технический университет \\
    имени Н.Э. Баумана \\
    (национальный исследовательский университет)>>\\
    (МГТУ им. Н.Э. Баумана)}\\
    \bottomrule[2pt]
    \bottomrule
\end{tabularx}

\noindent
\begin{tabularx}{\textwidth}{lX}
{\fontsize{12pt}{12pt}\selectfont ФАКУЛЬТЕТ}
    &
    {\fontsize{12pt}{12pt}\selectfont <<Информатика и системы управления>> (ИУ)}\\
    \hhline{~~}
    {\fontsize{12pt}{12pt}\selectfont КАФЕДРА} &
    {\fontsize{12pt}{12pt}\selectfont <<Информационная безопасность>> (ИУ8)}\\
    \hhline{~~}
\end{tabularx}
\vfill


\noindent
{\centering\rmfamily\bfseries\fontsize{20pt}{20pt}\selectfont
РАСЧЕТНО-ПОЯСНИТЕЛЬНАЯ ЗАПИСКА \\ \bigskip
\MakeTextUppercase{\RPZTitle} \\ \bigskip
НА ТЕМУ: \\ \bigskip
}


\vspace{1cm}

\noindent
\adjustbox{minipage=\textwidth}{\centering\bfseries\fontsize{20pt}{30pt}\selectfont \emph{\RPZTitleTheme}}

\vfill

\noindent
{\rmfamily\fontsize{12pt}{12pt}\selectfont Студент} \hfill $\stackrel[\text{\rmfamily\fontsize{9pt}{9pt}\selectfont (Группа)}]{}{\uline{\text{\rmfamily\fontsize{12pt}{12pt}\selectfont \RPZTitleGroup}}}$ \hfill\null \Signature{\RPZTitleStudent}


\noindent
{\rmfamily\fontsize{12pt}{12pt}\selectfont Научный руководитель} \hfill \Signature{\RPZTitleTeacher}

\vfill

{\centering\rmfamily\itshape\fontsize{14pt}{14pt}\selectfont \the\year~г. \par}
\newpage
\pagestyle{plain}

\Define{MAX}{максимум}
\Define{API}{test}

\newacronym{lvm}{LVM}{Logical Volume Manager}

\newacronym{test}{ТЕСТ}{Logical Volume Manager}

% Также можно использовать \Referat, как в оригинале
\begin{abstract}

    Отчет содержит \pageref{LastPage}\,стр.%
    \ifnum \totfig >0
    , \totfig~рис.%
    \fi
    \ifnum \tottab >0
    , \tottab~табл.%
    \fi
    %
    \ifnum \totbib >0
    , \totbib~источн.%
    \fi
    %
    \ifnum \totapp >0
    , \totapp~прил.%
    \else
    .%
    \fi


    Это пример каркаса расчётно-пояснительной записки, желательный к использованию в РПЗ проекта по курсу РСОИ
    \nocite{*}.

    Данный опус, как и более новые версии этого документа, можно взять по адресу (\url{https://github.com/latex-g7-32/latex-g7-32}).

    Текст в документе носит совершенно абстрактный характер.

\end{abstract}

%%% Local Variables: 
%%% mode: latex
%%% TeX-master: "rpz"
%%% End: 


\tableofcontents
\printglossary[title=ТЕРМИНЫ И ОПРЕДЕЛЕНИЯ, toctitle=ТЕРМИНЫ И ОПРЕДЕЛЕНИЯ,nonumberlist]
\printglossary[type=acronym, title=ПЕРЕЧЕНЬ СОКРАЩЕНИЙ И ОБОЗНАЧЕНИЙ, toctitle=ПЕРЕЧЕНЬ СОКРАЩЕНИЙ И ОБОЗНАЧЕНИЙ,nonumberlist]

\Introduction

Целью работы является создание всякой всячины. Для достижения поставленной цели необходимо решить следующие задачи:

\begin{itemize}
\item проанализировать существующую всячину;
\item спроектировать свою, новую всячину;
\item изготовить всякую всячину;
\item проверить её работоспособность.
\end{itemize}

Проверяем как у нас работают сокращения, обозначения и определения "---
\gls{MAX}
API 
\gls{API}
\gls{test}
с обратным прокси.





\chapter*{ОСНОВНАЯ ЧАСТЬ}
\addcontentsline{toc}{chapter}{ОСНОВНАЯ ЧАСТЬ}
\mainmatter

%\chapter{Технологический раздел}
\label{cha:impl}

В данном разделе описано изготовление и требование всячины. Кстати,
в Latex нужно эскейпить подчёркивание (писать <<\verb|some\_function|>> для \Code{some\_function}).

\ifPDFTeX
Для вставки кода есть пакет \Code{listings}. К сожалению, пакет \Code{listings} всё ещё
работает криво при появлении в листинге русских букв и кодировке исходников utf-8.
В данном примере он (увы) на лету конвертируется в koi-8 в ходе сборки pdf.

Есть альтернатива \Code{listingsutf8}, однако она работает лишь с
\Code{\textbackslash{}lstinputlisting}, но не с окружением \Code{\textbackslash{}lstlisting}

Вот так можно вставлять псевдокод (питоноподобный язык определен в \Code{listings.inc.tex}):

\begin{lstlisting}[style=pseudocode,caption={Алгоритм оценки дипломных работ}]
def EvaluateDiplomas():
    for each student in Masters:
        student.Mark := 5
    for each student in Engineers:
        if Good(student):
            student.Mark := 5
        else:
            student.Mark := 4
\end{lstlisting}

Еще в шаблоне определен псевдоязык для BNF:

\begin{lstlisting}[style=grammar,basicstyle=\small,caption={Грамматика}]
  ifstmt -> "if" "(" expression ")" stmt |
            "if" "(" expression ")" stmt1 "else" stmt2
  number -> digit digit*
\end{lstlisting}

В листинге~\ref{lst:sample01} работают русские буквы. Сильная магия. Однако, работает
только во включаемых файлах, прямо в \TeX{} нельзя.

% Обратите внимание, что включается не ../src/..., а inc/src/...
% В Makefile есть соответствующее правило для inc/src/*,
% которое копирует исходные файлы из ../src и конвертирует из UTF-8 в KOI8-R.
% Кстати, поэтому использовать можно только русские буквы и ASCII,
% весь остальной UTF-8 вроде CJK и египетских иероглифов -- нельзя.

\lstinputlisting[language=C,caption=Пример (\Code{test.c}),label=lst:sample01]{inc/src/test.c}

\else

Для вставки кода есть пакет \texttt{minted}. Он хорош всем кроме: необходимости Python (есть во всех нормальных (нет, Windows, я не про тебя) ОС) и Pygments и того, что нормально работает лишь в \XeLaTeX.

\ifdefined\NoMinted
Но к сожалению, у вас, по-видимому, не установлен Python или pygmentize.
\else
Можно пользоваться расширенным BFN:

\begin{listing}[H]
\begin{ebnfcode}
 letter = "A" | "B" | "C" | "D" | "E" | "F" | "G"
       | "H" | "I" | "J" | "K" | "L" | "M" | "N"
       | "O" | "P" | "Q" | "R" | "S" | "T" | "U"
       | "V" | "W" | "X" | "Y" | "Z" ;
digit = "0" | "1" | "2" | "3" | "4" | "5" | "6" | "7" | "8" | "9" ;
symbol = "[" | "]" | "{" | "}" | "(" | ")" | "<" | ">"
       | "'" | '"' | "=" | "|" | "." | "," | ";" ;
character = letter | digit | symbol | "_" ;
 
identifier = letter , { letter | digit | "_" } ;
terminal = "'" , character , { character } , "'" 
         | '"' , character , { character } , '"' ;
 
lhs = identifier ;
rhs = identifier
     | terminal
     | "[" , rhs , "]"
     | "{" , rhs , "}"
     | "(" , rhs , ")"
     | rhs , "|" , rhs
     | rhs , "," , rhs ;
 
rule = lhs , "=" , rhs , ";" ;
grammar = { rule } ;
\end{ebnfcode}
\caption{EBNF определённый через EBNF}
\label{lst:ebnf}
\end{listing}

А вот в листинге \ref{lst:c} на языке C работают русские комменты. Спасибо Pygments и Minted за это.

\begin{listing}[H]
\cfile{inc/src/test.c}
\caption{Пример — test.c} 
\end{listing}
\label{lst:c}

\fi
\fi
% Для вставки реального кода лучше использовать \texttt{\textbackslash lstinputlisting} (который понимает
% UTF8) и стили \Code{realcode} либо \Code{simplecode} (в зависимости от размера куска).




Можно также использовать окружение \Code{verbatim}, если \Code{listings} чем-то не
устраивает. Только следует помнить, что табы в нём <<съедаются>>. Существует так же команда \Code{\textbackslash{}verbatiminput} для вставки файла.

\begin{verbatim}
a_b = a + b; // русский комментарий
if (a_b > 0)
    a_b = 0;
\end{verbatim}

%%% Local Variables:
%%% mode: latex
%%% TeX-master: "rpz"
%%% End:

\section{Понятие идеала в кольце многочленов}
\label{cha:basedifinition}

В первой части научно-исследовательской работы даются базовые определения и понятия, которые будут использоваться на протяжении всей расчетно-пояснительной записки.
\\Пусть $R$ - коммутативное кольцо с единицей $1$.


\subsection{Понятие идеала}

Непустое подмножество $I$ кольца $R$ называется идеалом в $R$ (записывается $I \triangleleft R$), если:

1) для любых элементов $a,b \in I$ элемент $a + b \in I$;

2)для любых $a \in I, c \in R$ элемент $ac \in I$.

Идеал $I$ кольца $R$ называется главным, если существует такой элемент $a \in I$, что $ I = (a)$. Элемент $a$ называется порождающим( или образующим)  для идеала $I$.

Кольцо  $R$ называется кольцом главных идеалов, если каждый идеал кольца $R$ является главным.

\subsection{Базис идеала}

Обобщим по понятие главного идеала $I$ кольца $R$.  Пусть 
$a_{1}, ... , a_{k}$ - произвольные элементы кольца $R$.

Если множество 
$$ (a_{1}, ... , a_{k}) =\left\{ a_{1}r_{1} + ... + a_{k}r_{k};  r_{1}, ... , r_{k} \in R \right\} $$ есть идеал $I$ кольца $R$, тогда говорят, что элементы $a_{1}, ... , a_{k}$ составляют базис идеала $I$. Обозначается  $I = (a_{1}, ... , a_{k})$.

Можно заметить, что при определении базиса идеала нет требования на минимальное количество базисных элементов. Это связано с тем, что при добавлении к базису произвольного элемента идеала, мы получаем тот же самый базис.

Для любого элемента $a \in (a_{1}, ..., a_{k})$:
$$ (a_{1}, .. , a_{k}, a) = (a_{1}, ... ,a_{k}) $$ 

Легко убедиться в достоверности этого утверждения. По определению $ (a_{1}, ... , a_{k}, a) = \left\{ a_{1}r_{1} + ... + a_{k}r_{k} + ar;   r_{1}, ... , r_{k}, r \in R \right\} $. По определению идеала $I$ кольца $R$, если $a \in I, r \in R$, то элемент $ar \in I$. Тогда 
$ar = a_{1}r_{1}' + ... + a_{k}r_{k}';  r_{1}', ... , r_{k}', \in R $. Подставим это выражение в определение базиса $ (a_{1}, ... , a_{k}, a) = \left\{ a_{1}r_{1} + ... + a_{k}r_{k} + a_{1}r_{1}' + ... + a_{k}r_{k}';  r_{1}, r_{1}', ... , r_{k},  r_{k}' \in R \right\} = 
 \left\{ a_{1}(r_{1} + r_{1}') + ... + a_{k}(r_{k} + r_{k}');  r_{1}, r_{1}', ... , r_{k}, r_{k}' \in R \right\} $. Таким образом, при добавлении нового базисного элемента, который принадлежит идеалу, базис не изменится.
 
 \subsection{Идеал в кольцах многочленов}
 
 Пусть $K \left[ x_{1}, ... , x_{n} \right]$ - кольцо многочленов от переменных $x_{1}, ... ,x_{n}$ над полем $\mathcal{K}$.
 Рассмотрим идеалы в кольцах многочленов.
 
 Легко убедиться в том, что кольцо многочленов от нескольких переменных не является кольцом главных идеалов.
В кольце многочленов $K \left[ x_{1}, ... , x_{n} \right]$ выделим множество многочленов, у которых свободных член равен $0$. Данное множество образует идеал в этом кольце, обозначим его $I_{0}$. Пусть $ I_{0}  = (f), f \in K \left[ x_{1}, ... , x_{n} \right] $. Поскольку $x_{1} \in I_{0}$, значит $f$ - либо ненулевая константа (тогда $I_{0} =K \left[ x_{1}, ... , x_{n} \right] $, что является противоречием условия), либо $f = ax, a \in \mathcal{K}, a \neq 0$. Но $x_{2} \in I_{0}$, значит $f$ делит 
 $x_{2}$. Возникает противоречие.
 
Раз кольцо многочленов от нескольких переменных не является кольцом главных идеалов, значит имеет место понятие базиса идеала. В следующем параграфе будут подробно рассмотрены базисы идеала в кольце многочленов от нескольких переменных и выведены некоторые утверждения.







 
 
 


%%% Local Variables:
%%% mode: latex
%%% TeX-master: "rpz"
%%% End:

\section{Мономиальные базисы идеалов в кольцах многочленов} 

В данном разделе будут подробно рассмотрены мономиальные базисы идеалов в кольцах многочленов от нескольких переменных и введено понятие базиса Гребнера.

\subsection{Мономиальный порядок и его свойства}
Пусть $ M_{n} $ - множество мономов. 
Мономиальным упорядочением $ \prec $ на $ M_{n} $ называется линейный порядок, удовлетворяющий свойствам:

1)  $M \prec N \Rightarrow MP \prec NP,  \forall M, N, P \in M_{n}$ 

2)$ 1 \preceq M,  \forall M \in M_{n} $

Пусть $ M = x_{1}^{a_{1}}...x_{n}^{a_{n}} $, $ N = x_{1}^{b_{1}}...x_{n}^{b_{n}} $ - произвольные мономы из $ M_{n} $.
Приведем несколько примеров бинарных отношений на $M_{n}$, которые являются мономиальными упорядочениями:

1.  Лексикографическое упорядочение (lex):
$$ M \prec_{lex} N \Leftrightarrow (a_{1}, ... , a_{n}) \prec_{lex} (b_{1}, ..., b_{n}) $$

2. Сначала по степени, затем лексикографическое упорядочение (deglex):
$$ M \prec_{deglex} N \Leftrightarrow (deg(M), a_{1}, ... , a_{n}) \prec_{lex} (deg(N), b_{1}, ..., b_{n}) $$

3. Сначала по степени, затем обратное лексикографическое упорядочение (degrevlex):
$$ M \prec_{degrevlex} N \Leftrightarrow (deg(M), b_{n}, ... , b_{1}) \prec_{lex} (deg(N), a_{n}, ..., a_{1}) $$

Далее введем несколько обозначений, которые понадобятся нам в дальнейшем.

Пусть $ f= \sum_{a}p_{a}x^{a}$ - ненулевой полином в $K \left[ x_{1}, ... , x_{n} \right]$, и пусть $ \prec $ - мономиальное упорядочение.

1. Мультистепень полинома $f$ определяется так:
$$ multideg(f) = max(a \in Z_{\geq0}^{n}: p_{a} \neq 0) $$

2. Старший коэффициент полинома $f$:
$$ LC(f) = p_{multideg}(f) $$

3. Старший моном полинома $f$:
$$ LM(f) = x^{multideg}(f) $$ ( с коэффициентом 1)

4. Старший член полинома $f$:
$$ LT(f) = LC(f)LM(f) $$ 

\subsection{Мономиальные идеалы}

Введем понятие мономиального идеала в кольцое многочленов от нескольких переменных.

Идеал $I \triangleleft K \left[ x_{1}, ... , x_{n} \right]$ называется мономиальным, если существует подмножество $A \in Z_{\geq0}^{n}$ (которое может быть бесконечым), такое, что $I$ состоит из всех конечных сумм вида $\sum_{a \in A}h_{a}x^{a}$, где $h_{a} \in  K \left[ x_{1}, ... , x_{n} \right]$. Такой идеал будем обозначать $(x^{a}: a \in A)$.

\textbf{Лемма Диксона.} {
\it 
Любой мономиальный идеал $ I = ( x^{a} : a \in A)$ может быть представлен в виде $ I = (x^{a(1)}, ... , x^{a(s)}$, где $a(1), ..., a(s) \in A$. В частности, $I$ имеет конечный базис.
}

\subsection{Теорема Гильберта о базисе}
В предыдущем параграфе было  показано, что любой мономиальный идеал кольца многочленов от нескольких переменных иметь конечный базис. В этом разделе будет доказано, что любой идеал в кольце многочленов(не только мономиальный) имеет конечный базис.

\textbf{Теорема Гильберта о базисе.}{
\it
Каждый идеал $I \triangleleft K \left[ x_{1}, ... , x_{n} \right]$ является конечно порожденным, то есть $I = (g_{1}, ..., g_{s})$, где $g_{1}, ..., g_{s} \in I$.
}

\textbf{Доказательство.} 
Пусть $I \triangleleft K \left[ x_{1}, ... , x_{n} \right]$. Обозначим через $a_{t}$ множество элементов $a \in \mathcal{K}$, которое содержит в себе $LC(f)$, где $f \in I$ и $ multideg(f) = t $.
Легко видно, что $a_{i}$ является идеалом в $ \mathcal{K} $. (Если $a,b \in a_{i}$, то $ a \pm b \in a_{i}$, это легко увидеть, достаточно взять сумму или разность соответсвующих многочленов. Если $ c \in \mathcal{K}$, то $ca \in a_{i}$, это можно увидеть при умножении соответсвующего многочлена на $c$.)
Кроме того, имеем $$ a_{0} \subset a_{1} \subset a_{2} ... $$
другими словами последовательность идеалов $\left[ a_{i} \right]$ возрастающая. 
\\Пусть последовательность идеалов стабилизируется на $a_{r}$:
$$ a_{0} \subset a_{1} \subset a_{2} ... \subset a_{r} = a_{r+1} = ...$$
Пусть 
$$  a_{0}  = (a_{01}, a_{02}, ..., a_{0n_{0}})$$
$$  a_{r}  = (a_{r1}, a_{r2}, ..., a_{rn_{0}})$$
Для каждого $ i = 0, ...., r$ и $j = 0, ..., n_{i}$ пусть $f_{ij}$ - многочлен из $I$ степени $i$ со старшим коэффициентом $a_{ij}$. Покажем, что многочлены $f_{ij}$ образуют базис для $I$.
\\ Пусть $f$ - многочлен степени $d$ из $I$. Индукцией по $d$ докажем, что $f$ лежит в идеале, порожденном $f_{ij}$.
\\ Если $d > r$, то заметим, что старшие коэффициенты многочленов 
$$ X^{d-r}f_{r1}, ... , X^{d-r}f_{rn_{r}} $$
порождают $a_{d}$. Значит, существуют $ c_{1}, ... , c_{n_{r}} \in  \mathcal{K}$, такие, что многочлен
$$ f - c_{1}X^{d-r}f_{r1} - ... - c_{n_{r}}X^{d-r}f_{rn_{r}} $$
имеет степень $ < d$. Этот многочлен лежит в $I$.
Если $ d \leq r$, мы можем получить многочлен степени $ < d$, лежащий в $I$, вычитая некоторую линейную комбинацию
$$ f - c_{1}f_{d1} - ... - c_{n_{r}}f_{dn_{r}} $$
По индукции мы можем найти такой многочлен $g \in I$, который порожден $f_{ij}$ и $f - g = 0$, доказав таким образом теорему.

\subsection{Базис Гребнера}
В предыдущем параграфе была приведена и доказана фундаментальная теорема, связанная с базисом идеалов в кольцах многочленов.
Далее  будет введено понятие базиса Гребнера и описаны его особенности.

Пусть задано мономиальное упорядочение. Конечное подмножество $G = g_{1}, ... , g_{s}$ элементов идеала $I$ называется его базисом Гребнера, если 
$$ (LT(g_{1}), ..., LT(g_{s})) = ( LT(I) ) $$

Стоит отметить, что любой $I \triangleleft K \left[ x_{1}, ... , x_{n} \right]$ имеет базис Гребнера. Это легко увидеть из доказательства теоремы Гильберта о базисе.

Базис Гребнера является очень важным элементом в алгебраическом криптоанализе. Он помогает упрощать и решать алгебраические системы уравнений, выявлять различные особенности, связанные с той или иной системой алгебраических уравнений (например, существует эффективный критерий несовместимости системы). Поэтому, нахождение базиса Гребнера является одной из важных математических задач на данный момент. Далее будут рассмотрены три алгоритма нахождения базиса Гребнера (алгоритм Бухбергера, алгоритм $F4$ и алгоритм $F5$).

%%% Local Variables:
%%% mode: latex
%%% TeX-master: "rpz"
%%% End:

\backmatter %% Здесь заканчивается нумерованная часть документа и начинаются ссылки и
            
\Conclusion % заключение к отчёту

В результате проделанной работы стало ясно, что ничего не ясно...

%%% Local Variables: 
%%% mode: latex
%%% TeX-master: "rpz"
%%% End: 
%% заключение


% % Список литературы при помощи BibTeX
% Юзать так:
%
% pdflatex rpz
% bibtex rpz
% pdflatex rpz

\bibliographystyle{ugost2008}
\bibliography{rpz}

%%% Local Variables: 
%%% mode: latex
%%% TeX-master: "rpz"
%%% End: 



\appendix   % Тут идут приложения

\chapter{Картинки}
\label{cha:appendix1}

\blindtext
\begin{figure}
\centering
\caption{Картинка в приложении. Страшная и ужасная.}
\end{figure}

%%% Local Variables: 
%%% mode: latex
%%% TeX-master: "rpz"
%%% End: 

\chapter{Еще картинки}
\label{cha:appendix2}
\blindtext

\begin{figure}
\centering
\caption{Еще одна картинка, ничем не лучше предыдущей. Но надо же как-то заполнить место.}
\end{figure}

%%% Local Variables: 
%%% mode: latex
%%% TeX-master: "rpz"
%%% End: 


\end{document}